\section{Semantic XPath：树形结构化记忆的查询访问}
Semantic XPath\upcite{liu-etal-2026-semantic}聚焦结构化对话记忆的精准访问，结合树路径导航、位置筛选与连续语义评分，从层级记忆树中定位目标节点及子树，为下游回答与状态更新提供上下文。

\subsection{研究问题与方法定位}

在长期任务导向对话中，用户常围绕同一任务对象（如旅行计划、待办清单）反复查看、修改或回溯。其中一项具体活动不仅包含文本描述，还从属于特定日期、行程和历史版本。直接输入全量历史开销高昂；扁平检索虽能缩减输入，却易召回语义相似但所属日期或版本错误的条目。例如“在会议最密集的一天增加咖啡休息”，需要先按日期聚合评估各天活动，而非仅找一条含“会议”字样的文本。Semantic XPath据此将结构路径约束与软语义匹配结合，专门负责在树形记忆中锁定访问范围，再交由下游完成回答或节点增删改。

\subsection{树形数据模型与查询语法}

系统将记忆表示为有根树 $T=(V,E,r)$，非根节点具有唯一父节点，各节点包含类型标签与文本属性，树模式 $\Sigma$ 由数据自身的节点类型、属性与父子关系自然形成。旅行规划场景采用
\[
\text{Itinerary}\rightarrow\text{Version}\rightarrow\text{Day}\rightarrow\text{POI},
\]
依次表示行程、版本、日期和活动；上下层体现组成与归属关系，完整保留底层细节而非仅做有损摘要。

查询由步骤序列组成，单步写为 $S=A\,N\,[R]\,[P]$，其中路径轴 $A$ 与节点类型 $N$ 必选，位置选择器 $R$ 与语义谓词 $P$ 可选（表\ref{tab:sxpath-grammar}）。
\begin{table}[H]
\centering
\caption{Semantic XPath 查询步骤的组成}
\label{tab:sxpath-grammar}
\begin{tabularx}{\linewidth}{@{}p{20mm}p{43mm}X@{}}
\toprule
组成 & 示例写法 & 作用 \\
\midrule
路径 $A$ & \texttt{/}、\texttt{//} & 访问直接子节点或全部后代 \\
类型 $N$ & \texttt{Day}、\texttt{POI}、\texttt{*} & 按节点类型筛选或匹配任意类型 \\
位置 $R$ & \texttt{[i]}、\texttt{[-i]}、\texttt{[i:j]} & 按顺序、倒序或区间选取候选 \\
语义 $P$ & \texttt{[node\textasciitilde="workshop"]} & 计算候选与语义条件的相关程度 \\
\bottomrule
\end{tabularx}
\end{table}

例如，\texttt{//Day[3]/POI[node\textasciitilde="workshop"]}先定位第3天，再对该日期下的活动计算语义得分；存在多个行程或版本时，可通过前缀路径先限定目标范围。语义条件还支持通过 $1-P$ 取反，以筛选与指定概念无关的候选节点。

\subsection{语义评分与带权执行}

节点 $u$ 对语义谓词 $P$ 的满足程度定义为连续相关度 $\operatorname{Rel}(u,P)\in[0,1]$。除基于当前节点属性的局部评分外，查询还支持嵌套子查询，先对下层节点评分，再用平均（\texttt{avg}）、最小、最大或几何平均聚合下层证据；多个条件可通过乘积、均值及取反组合。

\begin{quote}\normalsize\ttfamily
//Day[avg(/POI[node\textasciitilde="conference"])]
\end{quote}
该查询通过聚合每日下属活动的语义得分来评价日期。示例中第2天的三个活动得分为$0.603$、$0.482$和$0.608$，均值约$0.565$，在候选日期中排名最高，从而被识别为会议最密集的一天。

执行器从根节点及初始权重1出发，维护带权集合 $W\subseteq V\times[0,1]$，按步骤执行路径扩展、类型筛选、位置选择与语义评分更新：
\begin{equation}
w_{\mathrm{new}}=w_{\mathrm{old}}\operatorname{Rel}(u,P).
\end{equation}
遇到聚合条件时递归求值下层子查询，最终按预设阈值或返回数量输出高权重节点及子树。

\subsection{范围约束示例}

以行程调整为例，假设当前版本第2天安排海滨公园，第5天安排山地徒步，第7天安排植物园和海港游船。面对“找出第7天下雨时需调整的户外活动”，纯向量检索易混入其他日期的户外活动；结构化访问则先限定第7天子树，再评估户外属性，准确锁定植物园与游船。单层日期约束虽可由带元数据过滤的扁平RAG实现，但Semantic XPath还能统一处理多级版本嵌套及依赖子节点聚合评分的复杂访问需求。

\subsection{实验设置与结果分析}

论文在旅行行程、待办清单和家庭配餐三个自建场景（每场景20个单轮、5个多轮对话）以及LoCoMo原生会话层级（Memory--Conversation--Session--Turn）上评测。回答模型采用GPT-5 mini与Gemini 3 Flash；相似度评分（XPath Sem）与Flat RAG使用Qwen3-Embedding-8B，蕴含评分（XPath Entail）使用\texttt{facebook/bart-large-mnli}。自建场景统计任务通过率（需检索正确数据且输出符合约束），LoCoMo统计模型裁判得分，并将XPath查询生成输入计入平均输入tokens。

\begin{table}[H]
\centering
\caption{三个自建场景的平均单轮结果}
\label{tab:sxpath-single}
\begin{tabular}{@{}llrr@{}}
\toprule
回答模型 & 方法 & 通过率（\%） & 平均tokens \\
\midrule
GPT-5 mini & In-context & 70.0 & 26\,216 \\
 & Flat RAG & 30.0 & 3\,250 \\
 & XPath Sem & 66.3 & 6\,166 \\
 & XPath Entail & 83.0 & 5\,691 \\
\midrule
Gemini 3 Flash & In-context & 79.0 & 25\,387 \\
 & Flat RAG & 37.7 & 3\,361 \\
 & XPath Sem & 64.3 & 6\,043 \\
 & XPath Entail & 70.7 & 5\,830 \\
\bottomrule
\end{tabular}
\end{table}

单轮实验（表\ref{tab:sxpath-single}）中，GPT-5 mini配合XPath Entail的通过率达83.0\%，较Flat RAG提升53.0个百分点，且token消耗仅为全量上下文的21.7\%；在Gemini 3 Flash上，XPath Entail（70.7\%）同样大幅优于Flat RAG（37.7\%），但略低于全量上下文（79.0\%）。

\begin{table}[H]
\centering
\caption{三个自建场景的平均多轮结果}
\label{tab:sxpath-multi}
\begin{tabular}{@{}llrr@{}}
\toprule
回答模型 & 方法 & 通过率（\%） & 平均tokens \\
\midrule
GPT-5 mini & In-context & 80.0 & 69\,097 \\
 & Flat RAG & 26.7 & 3\,327 \\
 & XPath Sem & 70.0 & 6\,294 \\
 & XPath Entail & 83.3 & 5\,850 \\
\midrule
Gemini 3 Flash & In-context & 56.7 & 76\,603 \\
 & Flat RAG & 26.7 & 3\,240 \\
 & XPath Sem & 76.7 & 6\,092 \\
 & XPath Entail & 80.0 & 5\,993 \\
\bottomrule
\end{tabular}
\end{table}

在包含多次版本修改的多轮实验（表\ref{tab:sxpath-multi}）中，XPath Entail在两种模型下均取得最高通过率（83.3\%与80.0\%），token用量仅为全量上下文的8.5\%与7.8\%；同时Entail稳定优于Sem，说明蕴含模型更适合判定条件满足关系。

\begin{table}[H]
\centering
\caption{LoCoMo结果（XPath输入含查询生成）}
\label{tab:sxpath-locomo}
\begin{tabular}{@{}llrr@{}}
\toprule
回答模型 & 方法 & 回答得分 & 平均输入tokens \\
\midrule
GPT-5 mini & In-context & 74.55 & 24\,857 \\
 & Flat RAG & 35.71 & 1\,341 \\
 & Semantic XPath & 65.19 & 3\,032 \\
\midrule
Gemini 3 Flash & In-context & 87.01 & 27\,499 \\
 & Flat RAG & 41.49 & 1\,996 \\
 & Semantic XPath & 73.25 & 3\,870 \\
\bottomrule
\end{tabular}
\end{table}

在LoCoMo（表\ref{tab:sxpath-locomo}）上，Semantic XPath仅用全量上下文12.2\%--14.1\%的输入tokens，便大幅超越Flat RAG，但仍低于全量上下文（如GPT-5 mini多跳题得分为43.26，低于全量上下文的58.87）。原文错误分析指出，自由对话缺乏显式业务树，跨分支多跳查询易在节点检索阶段出现漏召回或选错发言节点。

\subsection{方法局限与课设启发}

Semantic XPath兼具结构范围约束与评分可解释性，但依赖清晰的层级质量与准确的查询生成，处理跨树关联和复杂多跳仍有局限。与侧重写入期事件绑定的StructMem相比，Semantic XPath侧重读取期范围收缩；课程设计可将二者结合，在检索时先按会话、日期或版本收敛候选范围，再匹配具体事件条目。

% 原文定位：PDF第1--2页，图1--2；PDF第3页，\S3.1；PDF第3--4页，\S3.2；PDF第10页；PDF第4页，\S3.3、图2；PDF第5页，\S4.1；PDF第6--7页，\S4.3、图8。
